%% ============================================================================
%%  chapters/chapter5.tex
%%  فصل ۵: نتیجه‌گیری و پیشنهادات پژوهشی
%%  قالب رسمی پایان‌نامه و رساله دانشگاه خوارزمی
%% ============================================================================

\chapter{نتیجه‌گیری و پیشنهادات پژوهشی}
\label{chap:conclusion}

\section{مقدمه}
\label{sec:ch5_intro}

در این فصل پایانی، به جمع‌بندی جامع دستاوردها و یافته‌های پژوهش پرداخته می‌شود. این تحقیق با هدف حل چالش‌های مزمن دانشجویان و پژوهشگران دانشگاه خوارزمی در تدوین پایان‌نامه‌ها و رساله‌های علمی بر پایه سامانه قدرتمند لاتک و بسته بومی زی‌پرشین به انجام رسید. در ادامه، ابتدا خلاصه‌ای از فعالیت‌های انجام‌شده و نتایج کلیدی حاصل ارائه می‌گردد. سپس نوآوری‌ها و دستاوردهای علمی و فنی پژوهش فهرست شده، محدودیت‌های ذاتی کار تبیین می‌گردد و در نهایت، پیشنهادهایی روشن و کاربردی برای پژوهش‌های آتی و توسعه‌های بعدی قالب پیش روی علاقه‌مندان قرار می‌گیرد.

\section{جمع‌بندی پژوهش و یافته‌ها}
\label{sec:ch5_summary}

نگارش پایان‌نامه‌های دانشگاهی به زبان فارسی در بسترهای نرم‌افزاری سنتی، همواره به دلیل ناتوانی در مدیریت هم‌زمان جهت متن راست‌به‌چپ، تعبیه اصطلاحات و مراجع انگلیسی، فرمول‌بندی‌های چندسطحی ریاضی و عدم انطباق با مقررات دقیق شیوه‌نامه‌ها با مشکلات ساختاری مواجه بوده است. دانشگاه خوارزمی به عنوان یکی از قطب‌های علمی معتبر کشور، نیازمند قالبی یکپارچه، مدرن، ماژولار و دارای تضمین کیفی بود که بتواند تمامی سطوح تحصیلی (کارشناسی، کارشناسی ارشد، و دکتری تخصصی) را پوشش دهد.

در این پژوهش، با بررسی موشکافانه شیوه‌نامه مصوب سال ۱۴۰۲ دانشگاه خوارزمی و استخراج دقیق معیارهای ابعادی، کلاسی اختصاصی به نام \lr{\texttt{kharazmi-thesis.cls}} و بسته همراه \lr{\texttt{kharazmi.sty}} طراحی و پیاده‌سازی گردید. فرآیند توسعه شامل تنظیم دقیق هندسه صفحات (حاشیه‌های بالا ۳/۵، پایین ۲/۵، راست ۳/۰ و چپ ۳/۰ سانتی‌متر)، فاصله استاندارد ۲/۲۵ سانتی‌متر برای سرصفحه‌های متقارن، تعریف قلم رسمی بی زر با مقیاس ۱/۰۸ برای متن اصلی، قلم بی تیتر برای عناوین، و تایمز نیو رومن برای بخش‌های لاتین بود. همچنین با تفکیک ساختار سند به پوشه‌های مجزا، امکان نگارش ماژولار فصول و صفحات آغازین و پایانی بدون تداخل فراهم شد.

آزمون‌های جامع کامپایل چندمرحله‌ای نشان داد که سند حاصل فاقد هرگونه خطای مرگبار یا هشدار نویسه گم‌شده، دارای ارجاع‌دهی کامل دوزبانه در بیب‌تک با سبک \lr{\texttt{plain-fa}}، و برخوردار از فهرست‌های خودکار جدول‌ها، شکل‌ها، نمودارها و نمادها است.

\section{دستاوردها و نوآوری‌های پژوهش}
\label{sec:ch5_contributions}

مهم‌ترین دستاوردها و نوآوری‌های علمی و فنی این پروژه عبارتند از:
\begin{enumerate}
  \item \textbf{انطباق ۱۰۰ درصدی با شیوه‌نامه رسمی خوارزمی:} پیاده‌سازی تمام مؤلفه‌های ظاهری و فنی مطابق فایل رسمی ورد (\lr{\texttt{docx}}) و راهنمای ۵۲ صفحه‌ای دانشگاه خوارزمی، از جمله ساختار سرصفحه «عنوان \lr{|} صفحه» و حاشیه‌های اختصاصی.
  \item \textbf{رفع ریشه‌ای خطاهای نویسه‌های گم‌شده:} شناسایی و رفع اشکالات ناسازگاری یونیکد در قلم بی زر (نظیر همزه ترکیبی و کاراکتر ممیز فارسی) و اصلاح آنها با نویسه‌های استاندارد، به‌گونه‌ای که گزارش کامپایل دارای صفر هشدار نویسه گم‌شده است.
  \item \textbf{سامانه ماژولار و تفکیک‌شده:} ساختاردهی مستند در پوشه‌های مجزای \lr{\texttt{front/}}، \lr{\texttt{chapters/}} و \lr{\texttt{back/}}، که یادگیری، ویرایش و مدیریت متون طولانی را برای دانشجویان بسیار آسان می‌سازد.
  \item \textbf{مدیریت یکپارچه مراجع دوزبانه:} ایجاد الگوی استاندارد در پرونده \lr{\texttt{references.bib}} و استفاده از سبک کتاب‌شناسی \lr{\texttt{plain-fa}} که امکان استناد متقابل به منابع فارسی و انگلیسی را با جهت و نشانه‌گذاری صحیح تضمین می‌کند.
  \item \textbf{پشتیبانی جامع از عناصر تخصصی:} طراحی محیط‌های اختصاصی برای جدول‌های پیشرفته (\lr{\texttt{booktabs}})، شکل‌ها و زیرشکل‌ها (\lr{\texttt{subcaption}})، نمودارهای مستقل با فهرست مجزا (\lr{\texttt{chart}})، قطعه‌کدهای برنامه‌نویسی (\lr{\texttt{listings}})، و الگوریتم‌های فارسی (\lr{\texttt{algorithmic}}).
  \item \textbf{توسعه مستندات کامل و راهنمای کاربری:} نگارش دفترچه راهنمای جامع در پرونده \lr{\texttt{README.md}} به دو زبان فارسی و انگلیسی جهت راهنمایی دانشجویان در مراحل نصب، تنظیمات فونت و کامپایل.
\end{enumerate}

\section{محدودیت‌های پژوهش}
\label{sec:ch5_limitations}

با وجود دستیابی به تمامی اهداف تعیین‌شده، این پژوهش با محدودیت‌هایی روبه‌رو بوده است که از جمله آنها می‌توان به موارد زیر اشاره نمود:
\begin{itemize}
  \item \textbf{وابستگی به قلم‌های سیستم‌عامل:} قلم‌های سری بی (بی زر، بی تیتر، بی نازنین) ممکن است به صورت پیش‌فرض بر روی برخی توزیع‌های لینوکس نصب نباشند که نیازمند نصب دستی توسط کاربر است (البته قلم آزاد وزیرمتن به عنوان گزینه جایگزین در قالب پیش‌بینی شده است).
  \item \textbf{نیاز به توزیع‌های به‌روز تک:} استفاده از قابلیت‌های پیشرفته نیازمند توزیع‌های مدرن نظیر میک‌تک ۲۶ یا تک‌لایو ۲۰۲۳ به بعد است و در توزیع‌های بسیار قدیمی ممکن است نسخه‌های ناسازگار بسته‌ها مشکل‌ساز شوند.
  \item \textbf{محدودیت موتورهای بیب‌تک کلاسیک در ارقام:} سبک‌های قدیمی بیب‌تک در تبدیل خودکار برخی شماره‌های صفحات لاتین در مراجع نیازمند تنظیم فیلد زبانی \lr{\texttt{language = \{persian\}}} در فایل مراجع هستند.
\end{itemize}

\section{پیشنهادات برای پژوهش‌های آتی}
\label{sec:ch5_future_work}

در راستای ارتقا و گسترش هرچه بیشتر این سامانه، پیشنهادهای پژوهشی و توسعه‌ای زیر برای تحقیقات آتی توصیه می‌گردد:
\begin{enumerate}
  \item \textbf{طراحی رابط کاربری گرافیکی و افزونه وی‌اس‌کد:} توسعه یک افزونه اختصاصی برای محیط \lr{VS Code} یا سامانه آنلاین \lr{Overleaf} جهت خودکارسازی ایجاد پروژه جدید پایان‌نامه خوارزمی و تکمیل خودکار فیلدهای فراداده.
  \item \textbf{پشتیبانی از موتور بیب‌لاتک و بیبر:} بومی‌سازی کامل بسته \lr{\texttt{biblatex}} و پردازشگر \lr{\texttt{biber}} برای مدیریت مراجع پیچیده و استایل‌های متنوع ارجاع‌دهی (نظیر \lr{APA}، \lr{IEEE}، و \lr{Vancouver}) به زبان فارسی.
  \item \textbf{قالب‌بندی خودکار ارائه‌ها و اسلایدهای دفاع:} طراحی کلاس سازگار با بسته بیمر (\lr{\texttt{beamer}}) با نشان‌ها و رنگ‌بندی مصوب دانشگاه خوارزمی جهت برگزاری جلسات دفاع کارشناسی ارشد و دکتری.
  \item \textbf{یکپارچه‌سازی با ابزارهای هوش مصنوعی و مدل‌های زبانی:} اتصال قالب به ابزارهای خودکار اصلاح نگارش و فاصله‌گذاری مجازی فارسی با استفاده از مدل‌های زبانی بزرگ جهت ارتقای روانی متن پیش از ارسال نهایی.
\end{enumerate}

\section{کلام پایانی}
\label{sec:ch5_epilogue}

امید است قالب طراحی‌شده در این پایان‌نامه گامی مؤثر در جهت استانداردسازی انتشارات علمی دانشگاه خوارزمی، ارتقای جایگاه پژوهشی دانشگاه در مجامع ملی و بین‌المللی، و ایجاد تجربه‌ای لذت‌بخش و بی‌دغدغه برای دانشجویان و اساتید گرامی در فرآیند تدوین آثار ماندگار علمی باشد.
